Astronomers discovered an extraordinary binary system in the constellation of
Auriga during the course of the Apache Point Observatory Galactic Evolution
Experiment (APOGEE). In these questions, you will try to analyse the data and
recreate their discovery for yourself.

The research team is aiming to find compact stars in binary systems using the
radial velocity (RV) technique. They examined archival APOGEE spectra of
``single'' stars and measured the apparent variation of their RV within this
data. Among ${\sim}200$ stars with the highest accelerations, researchers
searched for periodic photometric variations in data from the All-Sky
Automated Survey for Supernovae (ASAS-SN) that might be indicative of
transits, ellipsoidal variations or starspots. After this process, they
spotted a star named 2M05215658+4359220, with a large variation in RV and
photometric variability.

2.1. The following table presents the radial velocity measurements of
2M05215658+4359220 during three epochs of APOGEE spectroscopic observation.
Here we assume the variation of its RV is due to the existence of an unseen
companion. The proper motion of the stars can be ignored.

\begin{center}
Table 3. APOGEE Radial Velocity Measurements of 2M05215658+4359220

\medskip
\begin{tabular}{cccc}
Observation & MJD & RV (km\,s$^{-1}$) & Uncertainty (km\,s$^{-1}$) \\
1 & 56204.9537 & $-37.417$ & 0.011 \\
2 & 56229.9213 & 34.846 & 0.010 \\
3 & 56233.8732 & 42.567 & 0.010 \\
\end{tabular}
\end{center}

\begin{description}
\item[(D2.1.1)] (6 points) Use the data and a simple linear model to obtain
  an initial estimate of the apparent \textbf{maximum acceleration} of the
  star:
  \[ a_{\max} = \left.\frac{\Delta RV}{\Delta t}\right|_{\max}, \quad
     \mbox{unit: } \mathrm{km\,s^{-1}\,day^{-1}} \]
\item[(D2.1.2)] (9 points) Now use the data to obtain an initial estimate of
  the \textbf{mass} of its unseen companion.
\end{description}

2.2. After discovering this peculiar star, astronomers conducted follow-up
observations using the 1.5-m Tillinghast Reflector Echelle Spectrograph
(TRES) at the Fred Lawrence Whipple Observatory (FLWO) located on Mt.\
Hopkins in Arizona, USA. The following table presents the RV measurements
using this instrument:

\begin{center}
Table 4. TRES Radial Velocity Measurements of 2M05215658+4359220

\medskip
\begin{tabular}{ccc}
MJD & RV (km/s) & Uncertainty (km/s) \\
58006.9760 & 0 & 0.075 \\
58023.9823 & $-43.313$ & 0.075 \\
58039.9004 & $-27.963$ & 0.045 \\
58051.9851 & 10.928 & 0.118 \\
58070.9964 & 43.782 & 0.075 \\
58099.8073 & $-30.033$ & 0.054 \\
58106.9178 & $-42.872$ & 0.135 \\
58112.8188 & $-44.863$ & 0.088 \\
58123.7971 & $-25.810$ & 0.115 \\
58136.6004 & 15.691 & 0.146 \\
58143.7844 & 34.281 & 0.087 \\
\end{tabular}
\end{center}

\begin{description}
\item[(D2.2.1)] (14 points) \textbf{Plot} the diagram of RV variation
  (measured with TRES) versus time on your graph paper and label it as
  Figure 4. Draw a suitable \textbf{sinusoidal function} to fit the given
  data. \textbf{Estimate} the orbital period ($P_{\mathrm{orb}}$) and radial
  velocity semi-amplitude ($K$) from your plot.
\item[(D2.2.2)] (4 points) If the star is moving in a circular orbit,
  calculate the \textbf{minimum value} of the orbital radius
  ($r_{\mathrm{orb}}$) of the star in units of both $R_\odot$ and au.
\item[(D2.2.3)] (7 points) The mass function of a binary system is defined
  as:
  \[ f(M_1, M_2) = \frac{(M_2 \sin i_{\mathrm{orb}})^3}{(M_1 + M_2)^2} \]
  where the subscript ``1'' represents the primary star and ``2'' represents
  its companion. The parameter $i_{\mathrm{orb}}$ is the orbital inclination
  of the binary system. This mass function can also be expressed in terms of
  observable parameters. Calculate the \textbf{mass function of this system}
  in units of $M_\odot$.
\end{description}

2.3. Based on a detailed analysis of APOGEE, TRES spectra and GAIA parallax
measurements, astronomers derived the following stellar parameters:

\begin{center}
Table 5. Selected Physical Properties of 2M05215658+4359220

\medskip
\begin{tabular}{ccccc}
Effective & Surface Gravity & Parallax & Measured Rotation & Bolometric \\
Temperature & $\log g$ (cm\,s$^{-2}$) & $\pi$ (mas) & Velocity & Flux \\
$T_{\mathrm{eff}}$ (K) & & &
  $v_{\mathrm{rot}}\sin i$ (km\,s$^{-1}$) &
  $F$ (J\,s$^{-1}$\,m$^{-2}$) \\
$4890 \pm 130$ & $2.2 \pm 0.1$ & $0.272 \pm 0.049$ & $14.1 \pm 0.6$ &
  $(1.1 \pm 0.1) \times 10^{-12}$ \\
\end{tabular}
\end{center}

Photometric observations indicate that the period of its light curve is
identical to its orbital period, thus we may assume that the rotation period
satisfies $P_{\mathrm{rot}} = P_{\mathrm{orb}} \equiv P$, and the inclination
satisfies $i_{\mathrm{orb}} = i_{\mathrm{rot}} \equiv i$.

\begin{description}
\item[(D2.3.1)] (16 points) \textbf{Calculate} the luminosity ($L_1$, in
  unit of $L_\odot$), radius ($R_1$, in unit of $R_\odot$), sine of the
  inclination angle ($\sin i$), as well as mass ($M_1$, in unit of
  $M_\odot$) of the visible star. Please \textbf{include} the uncertainty in
  your results.
\item[(D2.3.2)] (4 points) \textbf{Choose} the correct type of this star
  from the following options: (1) Blue Giant (2) Yellow main sequence star
  (3) Red Giant (4) Red main sequence star (5) White Dwarf.
\item[(D2.3.3)] (10 points) Based on the mass function $f(M_1, M_2)$ of the
  binary system, \textbf{plot} the rough relationship of $M_2$ (as vertical
  axis) and $M_1$ (as horizontal axis) on your graph paper and label it as
  Figure 5. Plot the most probable relation (by using $\sin i$), upper limit
  (with $\sin i + \Delta\sin i$) and lower limit (with
  $\sin i - \Delta\sin i$) derived in (D2.3.1).
\item[(D2.3.4)] (5 points) Draw a vertical \textbf{shadowed region} of
  $[M_1 - \Delta M_1, M_1 + \Delta M_1]$, as well as two horizontal
  \textbf{dashed lines} showing the maximum mass of the white dwarf and
  neutron star, on your Figure 5. What is the possible mass of the invisible
  companion, and what kind of celestial object could it be?
\end{description}
